\chapter{Metodologi}
\label{chap:metodologi}

% Ubah bagian-bagian berikut dengan isi dari metodologi

Penelitian ini dilaksanakan sesuai dengan metodologi yang terstruktur dan sistematis untuk memecahkan masalah deteksi konjungtivitis pada citra mata yang memiliki derau (\textit{noise}) dan ketidakseimbangan kelas.

\section{Alur Penelitian}
\label{sec:alur-penelitian}

Alur pengerjaan penelitian ini dirancang melalui tahapan terintegrasi yang digambarkan sebagai berikut:
\begin{enumerate}
  \item \textbf{Pengumpulan Dataset}: Mengambil dataset publik konjungtivitis dari Kaggle.
  \item \textbf{Pembagian Dataset}: Memisahkan data menjadi subset latih dan uji secara proporsional.
  \item \textbf{Prapemrosesan Citra}: Melakukan segmentasi sklera dan memperjelas pembuluh darah dengan CLAHE.
  \item \textbf{Augmentasi Data}: Menyeimbangkan sebaran kelas minoritas pada data latih.
  \item \textbf{Perancangan dan Pelatihan Model}: Membangun arsitektur \textit{Attention-driven CNN} berbasis ResNet50.
  \item \textbf{Evaluasi Model}: Menganalisis hasil pengujian dengan metrik klasifikasi data tidak seimbang.
\end{enumerate}

\section{Pengumpulan dan Pembagian Dataset}
\label{sec:pengumpulan-dataset}

Dataset citra mata untuk deteksi konjungtivitis diperoleh secara publik dari platform Kaggle. Dataset ini memiliki karakteristik tidak seimbang (\textit{imbalanced}), dengan distribusi total sebanyak 1.357 citra yang terdiri dari:
\begin{itemize}
  \item Kelas \textbf{Normal}: 986 citra.
  \item Kelas \textbf{Konjungtivitis}: 371 citra.
\end{itemize}

Untuk melatih dan menguji model secara adil, dataset dibagi menjadi data latih (\textit{training set}) sebesar 80\% dan data uji (\textit{testing set}) sebesar 20\% menggunakan metode \textit{Stratified Split}. Penggunaan \textit{Stratified Split} menjamin bahwa rasio kelas Normal (72,6\%) dan kelas Konjungtivitis (27,4\%) tetap sama baik pada data latih maupun data uji, sehingga mencegah bias evaluasi.

\section{Prapemrosesan dan Augmentasi Data}
\label{sec:prapemrosesan-dan-augmentasi}

\subsection{Segmentasi Sklera dan Fitur Vaskular}
\label{subsec:segmentasi-sklera-vaskular}

Proses prapemrosesan citra dilakukan untuk menormalisasi pencahayaan, menghilangkan pantulan cahaya, mereduksi derau non-okular, serta mempertegas tanda-tanda inflamasi. Tahapan pemrosesan adalah sebagai berikut:
\begin{enumerate}
  \item \textbf{Normalisasi Histogram}: Menyeragamkan tingkat kecerahan dan kontras grayscale pada seluruh citra mata untuk mengompensasi variasi pencahayaan luar ruangan.
  \item \textbf{Restorasi Glare (Pantulan Cahaya)}: Mengonversi citra ke ruang warna CIE LAB. Piksel dengan nilai kecerahan tinggi ($L > 220$) pada kanal L dideteksi sebagai \textit{glare} akibat air mata atau cahaya eksternal. Area ini disamarkan (\textit{inpainting}) menggunakan informasi piksel tetangga.
  \item \textbf{Segmentasi Sklera}: Area putih mata (sklera) diisolasi dengan mengonversi citra ke ruang warna HSV dan menerapkan ambang batas (\textit{color thresholding}) warna putih sklera. Derau kelopak mata, bulu mata, dan pupil dibuang menggunakan deteksi kontur dan operasi morfologi (\textit{opening} dan \textit{closing}), dilanjutkan dengan pemotongan (\textit{cropping}) area sklera.
  \item \textbf{Ekstraksi Fitur Vaskular}: Citra sklera yang telah terpotong dikonversi untuk memperkuat kanal merah pembuluh darah melalui formulasi indeks kemerahan. Selanjutnya, algoritma \textit{Contrast Limited Adaptive Histogram Equalization} (CLAHE) diterapkan untuk menonjolkan urat-urat merah pembuluh darah yang mengalami pembengkakan dengan meratakan kontras lokal secara adaptif.
\end{enumerate}

\subsection{Augmentasi Data Latih}
\label{subsec:augmentasi-data-latih}

Untuk mengatasi bias kelas mayoritas (Normal) akibat ketidakseimbangan data, dilakukan teknik augmentasi data (\textit{data augmentation}) secara dinamis khusus pada kelas Konjungtivitis di subset latih. Augmentasi yang diterapkan meliputi:
\begin{itemize}
  \item Rotasi acak (\textit{random rotation}) dalam rentang $\pm 15$ derajat.
  \item Pembalikan horizontal (\textit{horizontal flip}) secara acak.
  \item Pengaturan kecerahan (\textit{brightness adjustment}) secara acak antara 0.8 hingga 1.2.
  \item Pergeseran lebar dan tinggi (\textit{width and height shift}) sebesar 10\%.
\end{itemize}
Dengan augmentasi ini, jumlah citra kelas minoritas seolah-olah seimbang selama proses pelatihan.

\section{Perancangan Model Klasifikasi}
\label{sec:perancangan-model}

Arsitektur model yang dirancang adalah \textit{Attention-driven CNN} dengan menggunakan ResNet50 yang telah terlatih pada ImageNet sebagai \textit{backbone} ekstraktor fitur dasar. Peta fitur (\textit{feature maps}) yang dihasilkan oleh blok konvolusi terakhir ResNet50 dialirkan ke dalam blok \textit{Spatial Attention Mechanism} untuk menghitung peta perhatian spasial. Blok perhatian ini memberikan bobot piksel yang lebih tinggi pada area yang menunjukkan vaskularisasi/urat merah, dan menekan area sklera yang bersih. Struktur implementasi model diilustrasikan pada Kode Program \ref{lst:attentioncnn}.

\lstinputlisting[
  language=Python,
  caption={Implementasi model Attention-driven CNN berbasis ResNet50 dan Spatial Attention.},
  label={lst:attentioncnn}
]{program/attention-cnn.py}

\section{Evaluasi Model dan Skenario Eksperimen}
\label{sec:evaluasi-model}

Karena dataset yang digunakan bersifat \textit{imbalanced}, evaluasi kinerja model usulan wajib menyertakan beberapa metrik evaluasi selain akurasi global (\textit{accuracy}) guna menghindari interpretasi bias. Metrik yang digunakan adalah:
\begin{enumerate}
  \item \textbf{Confusion Matrix}: Menyediakan visualisasi frekuensi prediksi \textit{True Positive} (TP), \textit{True Negative} (TN), \textit{False Positive} (FP), dan \textit{False Negative} (FN).
  \item \textbf{Precision}: Mengukur ketepatan prediksi positif, dirumuskan sebagai $TP / (TP + FP)$.
  \item \textbf{Recall (Sensitivity)}: Mengukur kemampuan mendeteksi kelas positif (Konjungtivitis), dirumuskan sebagai $TP / (TP + FN)$.
  \item \textbf{F1-Score}: Rata-rata harmonik dari Precision dan Recall yang sangat representatif untuk data tidak seimbang, dirumuskan sebagai $2 \times \frac{Precision \times Recall}{Precision + Recall}$.
  \item \textbf{AUC-ROC}: Area di bawah kurva karakteristik operasi penerima untuk menilai kemampuan pemisahan kelas.
\end{enumerate}

Skenario eksperimen dirancang dengan membandingkan model usulan terhadap beberapa baseline kontrol untuk membuktikan efektivitas masing-masing komponen. Rangkuman skenario eksperimen disajikan dalam Tabel \ref{tab:skenario-baseline}.

\begin{table}[!ht]
\centering
\footnotesize
\caption{Rangkuman Skenario Eksperimen dan Baseline}
\label{tab:skenario-baseline}
\begin{tabular}{|p{3cm}|p{4.5cm}|p{4.5cm}|}
\hline
\textbf{Parameter Evaluasi} & \textbf{Model Baseline (Kontrol)} & \textbf{Sistem Usulan} \\
\hline
\textit{Prapemrosesan Citra} & Citra mata utuh tanpa pemotongan (\textit{global classification}) dengan derau kelopak/bulu mata. & Segmentasi area sklera secara otomatis dilanjutkan dengan \textit{Vessel Enhancement} (CLAHE). \\
\hline
\textit{Arsitektur CNN} & Model CNN standar tanpa mekanisme perhatian (\textit{attention}). & Model CNN dengan integrasi \textit{Spatial Attention Block} pada ujung \textit{backbone} ResNet50. \\
\hline
\textit{Metrik Evaluasi Utama} & Akurasi global saja. & \textit{F1-Score}, \textit{Precision}, \textit{Recall}, \textit{Confusion Matrix}, dan kurva AUC-ROC. \\
\hline
\textit{Penanganan Imbalance} & Tanpa penanganan ketidakseimbangan data. & Augmentasi khusus kelas Konjungtivitis dan penerapan \textit{Class Weighting} pada \textit{loss function}. \\
\hline
\end{tabular}
\end{table}

